



















































































































































\begin{definition}[(Approximate) Kernel Density Estimation $\mathsf{KDE}_f(n, m, \eps, B)$]
Let $f: [0, B] \rightarrow [0, 1]$ be a real function and define the (kernel) function $k(x, y) = f(\|x-y\|_2^2)$.

Suppose we are given as input $2n$ points $x^{(1)}, \cdots, x^{(n)}, y^{(1)}, \cdots, y^{(n)} \in \R^m$ with the guarantee that $\|x^{(i)} - y^{(j)}\|_2^2 \le B$ for all $i, j \in [n]$.
Define the kernel matrix $K \in [0,1]^{n \times n}$ by $K[i, j] = f(\|x^{(i)}- y^{(j)}\|_2^2)$ for $i, j \in [n]$.
Then given additionally a vector $u \in \R^n$, one needs to output a vector $v \in \R^n$ such that
$\|v - K u\|_{\infty} \le \eps\|u\|_{1}$.
\end{definition}







































































































































































































































\section{Preliminaries}






























\subsection{SETH and known hard problems}
\label{sec:hardproblems}





















We now introduce several variants of the Bichromatic Closest Pair problem. 
\begin{problem}[Hamming (Exact) Bichromatic Closest Pair]
\label{prob:HBCP}
$\emph{\textsf{Hamming-BCP}}(n, m)$: Given two sets $A = \{x^{(1)}, \cdots, x^{(n)}\}, B = \{y^{(1)}, \cdots, y^{(n)}\} \subset \{0, 1\}^m$, compute $\min_{i, j \in [n]} \|x^{(i)} - y^{(j)}\|_2^2$.
\end{problem}
\begin{theorem}[\cite{AW15}]
Assuming SETH, for every $q \in (0, 1)$, there exists $C > 0$ such that if $m > C\log n$, then $\mathsf{Hamming}$-$\mathsf{BCP}(n, m)$ cannot be solved in time $O(n^{2-q})$ for any constant $q>0$.
\end{theorem}

Similarly, one can define the Hamming approximate BCP problem and its decision version.
\begin{problem}[Hamming Approximate BCP]
\label{prob:HABCP}
$\emph{\textsf{Hamming-Apx-BCP}}(n, m, \mu)$:
Given two sets $A, B$ as in Problem \ref{prob:HBCP} as well as $\mu \in \R_+$, output $d \in \R$ such that $\min_{i, j \in [n]} \|x^{(i)} - y^{(j)}\|_2^2 \le d \le (1 + \mu) \min_{i, j \in [n]} \|x^{(i)} - y^{(j)}\|_2^2$.
\end{problem}
\begin{theorem}[\cite{Rub18}] \label{bcp1}
Assuming SETH, for every $q > 0$, there exist $C > 0, \mu > 0$ such that if $m > C\log n$, then $\emph{\textsf{Hamming-Apx-BCP}}(n, m, \mu)$ cannot be solved in time $O(n^{2-q})$ for any constant $q>0$.
\end{theorem}

In the low dimensional regime $m = o(\log n)$, we consider the hardness of the $\ell_2$ BCP problem.
\begin{problem}[$\ell_2$ (Exact) Bichromatic Closest Pair]
$\ell_2\emph{\textsf{-BCP}}(n, m, D)$: Given two sets $A = \{x^{(1)}, \cdots, x^{(n)}\}, B = \{y^{(1)}, \cdots, y^{(n)}\} \in \Z^m$ such that $\max_{i, j \in [n]} \|x^{(i)} - y^{(j)}\|_2^2 \le D$, compute $\min_{i, j \in [n]} \|x^{(i)} - y^{(j)}\|_2^2$.
\end{problem}
\begin{theorem}[\cite{ChenMinIP}] \label{bcp2}
Assuming SETH, for every $q > 0$, there exists $C_1, C_2 > 0$ such that if $m > C_1^{\log^* n}$ and $D > m^{C_2^{\log^* n} \cdot (\log n)/m}$, then $\ell_2$\emph{\textsf{-BCP}}$(n, m, D)$ cannot be solved in time $O(n^{2-q})$ for any constant $q>0$.
\end{theorem}

To unify the hardness results for BCP in different dimension regimes, we view the Hamming BCP problem 
as an $\ell_2$ BCP problem with $D = m$. 
Formally, we combine Theorem~\ref{bcp1} and Theorem~\ref{bcp2} as follows.
\begin{theorem} \label{thm:mainbcp}
Assuming SETH, for every $q > 0$, there exists $C > 0$ such that $\ell_2$\emph{\textsf{-BCP}}$(n, m, D)$
cannot be solved in time $O(n^{2-q})$ for any constant $q>0$ if either of the following holds: (1) $m > C \log n, D = m$, or (2) $m > C^{\log^* n}, D > m^{C^{\log^* n} \cdot (\log n)/m}$.
\end{theorem}

\subsection{Kernels of interest}
In this work we focus primarily on three kernels $k(x, y) = f(\|x - y\|_2^2)$:
\begin{itemize}
\item Gaussian kernel $f(x) = e^{-x}$;
\item Rational Quadratic kernel $f(x) = 1/(1+x)^\sigma$ for $\sigma \ge 1$ a parameter;
\item $t$-Student kernel $f(x) = 1/(1+x^{\rho})$ for $\rho \ge 1$ a parameter.\footnote{$t$ in the name of the kernel in principle should be the name of the parameter. We here use $\rho$ as parameter while keeping the name \emph{$t$-Student kernel} unchanged.}
\end{itemize}
Rational Quadratic kernel and $t$-Student kernel are two typical kernels with mild decay (as opposed to the rapid decay of Gaussian kernel). This property is abstracted by Backurs et al. \cite{BCIS18} in the definition of a smooth kernel. Here we only focus on decreasing radial kernels.
\begin{definition}[smooth kernel]
A decreasing radial kernel $k(x, y) = f(\|x - y\|_2^2)$ is $(L, t)$ smooth if for any $0 < a < b$,
\[\frac{f(a)}{f(b)} \le L \left(\frac ba \right)^t.\]
\end{definition}
By calculation one can verify that Rational Quadratic kernel and $t$-Student kernel are respectively $(1, 1)$- and $(1, \rho)$-smooth.

\paragraph{Positive definite kernels}
Most commonly-studied kernels are positive definite. We will find that our lower bound approach takes a particularly nice form for such kernels.

\begin{definition}[Positive definite kernel]
A kernel $k: \R^m \times \R^m \rightarrow \R$ is positive definite if for any $n$ points $x_1, \cdots, x_n \in \R^m$, the Gram matrix $G = [k(x_i, x_j)]_{i, j \in [n]}$ is always positive definite.
\end{definition}
For radial kernels, we have the following concise characterization of positive definite kernels.
\begin{definition}
Let $f:\R_{\ge 0} \rightarrow \R$ be a real function. We say $f$ is absolutely monotone if $G \in C^{\infty}(\R_{\ge 0})$ and $f^{(k)}(t)$ for all $k \in \N$ and $t \ge 0$, and we say $f$ is completely monotone if $G \in C^{\infty}(\R_{\ge 0})$ and $(-1)^{k} \cdot f^{(k)}(t) \ge 0$ for all $k \in \N$ and $t \ge 0$.
\end{definition}
\begin{theorem}[Schoenberg's characterization] \label{schoenberg} Let $f:\R_{\ge 0} \rightarrow \R$ be a real function. Then the kernel $k(x, y) = f(\|x-y\|_2^2)$ is positive definite if and only if $f$ is completely monotone on $\R_{\ge 0}$.
\end{theorem}

\subsection{Tools from linear algebra}
\label{subsec:CauchyBinetTool}

As in prior work, we will make use of the Cauchy-Binet formula.

\begin{lemma}[Cauchy-Binet formula]
Let $k>0$ be an integer, and for functions $A: [k] \times \N \rightarrow \R$ and $B: \N \times [k] \rightarrow \R$, define the matrix $C \in \R^{k \times k}$ by, for $i, j \in [k]$, $C[i, j] = \sum_{\ell=0}^{\infty} A[i, \ell]B[\ell, j]$.
If the sum converges absolutely for all $i, j$, then
\[\det(C) = \sum_{1 \le \ell_1 < \cdots < \ell_k}
\det(A[\ell_1, \cdots, \ell_k]) \cdot \det(B[\ell_1, \cdots, \ell_k]).\]
Here $A[\ell_1, \cdots, \ell_k]$(resp. $B[\ell_1, \cdots, \ell_k]$) is the $k \times k$ matrix obtained from $A$ (resp. $B$) by taking the columns (resp. rows) $\ell_1, \cdots, \ell_k$.
\end{lemma}

\section{Main reduction from BCP}
\label{sec:mainreduction}
Most of our new lower bounds are based on the reduction below from Hamming or $\ell_2$ Exact Bichromatic Closest Pair to KDE. This reduction generalizes a framework developed in \cite{ACSS20} to accommodate different parameter regimes. In this section we will give the outline of the reduction and establish an upper bound on $1/\eps$ determined by the quantity $\tau(M)$.

\begin{definition}
\label{def:countingmatrix}
Let $D > 0$ be an integer. Fix vectors $\alpha, \beta \in \R^D$ and $c \in \R$. Suppose $c + \alpha_{\ell}\beta_{p} \in [0, 1]$ for all $\ell, p \in [D]$. Then for function $f: [0, 1] \rightarrow \R$, the counting matrix $M = M(D;f, \alpha, \beta)$ is a $D \times D$ matrix defined by $M[\ell, p] = f(c + \alpha_{\ell}\beta_p), \ell, p \in [D]$.
\end{definition}
\begin{theorem}
Let $\alpha \in \R^D$ be a fixed vector and $\beta \in \R^D$ be the identity vector defined by $\beta_p = p$. If $\KDE_f(n,m,\eps)$ can be solved in $T(n, m, \eps)$ time, then $\ell_2$\emph{\textsf{-BCP}}$(n, m, D)$ with $m = n^{o(1)}, D = n^{o(1)}$ can be solved in $T(n, m+1, (3n\tau(M))^{-1}) \cdot n^{o(1)} + n^{1 + o(1)}$ time, where $M$ is the $D \times D$ counting matrix associated with $f, \alpha, \beta$ and
\[\tau(M) = \max_{0 \ne b \in \R^{D}} \frac{\|M^{-1}b\|_{\infty}}{\|b\|_{\infty}}.\]
\end{theorem}

\begin{proof}
Given two sets $X = \{x^{(1)}, \cdots, x^{n}\}, Y = \{y^{(1)}, \cdots, y^{(n)}\} \subseteq \Z^m$ such that $\max_{i, j \in [n]} \|x^{(i)} - y^{(j)}\|_2^2 \le D$, let $W \in \N^{D \times n}$ denote the \emph{distance count matrix} defined by $W[p, i] = \#[j \in [n]: \|x^{(i)} - y^{(j)}\|_2^2 = p]$. Then the matrix product $U = M \times W$ gives
\begin{align}
U[\ell, i] &= \sum_{p=1}^D M[\ell, p] \cdot W[p, i] \notag\\ 
&=\sum_{p=1}^D f\big(c + \alpha_\ell\beta_p\big) 
\sum_{j \in [n]} \ind\left[\|x^{(i)} - y^{(j)}\|_2^2 = p\right] \notag\\
&=\sum_{j \in [n]} f\left(c + \alpha_\ell \cdot \beta\left[\|x^{(i)} - y^{(j)}\|_2^2\right]\right). \label{eq:kernelsum} 
\end{align}
We note that when $\beta$ is the identity vector, the summation (\ref{eq:kernelsum}) can be formulated as a KDE instance.
More specifically, let $\tilde x_{\ell}^{(i)}, \tilde y_{\ell}^{(j)} \in \R^{m+1}$ be defined by $\tilde x_{\ell}^{(i)}[k] = \sqrt{\alpha(\ell)} x^{(i)}[k]$ if $k \in [m]$ and $\tilde x_{\ell}^{(i)}[m+1] = \sqrt c$;
$\tilde y_{\ell}^{(j)}[k] = \sqrt{\alpha(\ell)} y^{(j)}[k]$ if $k \in [m]$ and $\tilde y_{\ell}^{(j)}[m+1] = 0$. Then
\[U[\ell, i] = \sum_{j \in [n]} f(\|\tilde x_{\ell}^{(i)} - \tilde y_{\ell}^{(j)}\|_2^2).\]
Therefore we have the following algorithm for $\ell_2$\textsf{-BCP}$(n, m, D)$ using $\KDE_f$ as a subroutine. 
On input $X = \{x^{(1)}, \cdots, x^{(n)}\}, Y = \{y^{(1)}, \cdots, y^{(n)}\}$ such that $\max_{i, j \in [n]} \|x^{(i)} - y^{(j)}\|_2^2 \le D$:
\begin{enumerate}
\item For $\ell \in [D]$, construct vectors $\tilde x_{\ell}^{(i)}, \tilde y^{(j)}_{\ell}, i, j \in [n]$.
Then approximate the $\ell$-th row of $U$: $\hat U[\ell] \approx U[\ell] = K_{\ell} \times \bo$ using the $\KDE_f$ oracle.
\item Compute $\hat W = M^{-1} \times \hat U$, and round each entry to the closest integer.
\end{enumerate}

We claim that if we call the $\KDE_f$ subroutine with $\eps = (3n\tau(M))^{-1}$, then the distance count matrix $W$ is exactly recovered after the rounding step.
Indeed, for fixed $i \in [n]$, letting $W[\cdot, i], U[\cdot, i]$ respectively denote the $i$-th column of matrix $W$ and $U$, we have
\[\|\hat W[\cdot, i] - W[\cdot, i]\|_{\infty}
=\|M^{-1}(\hat U[\cdot, i] - U[\cdot, i])\|_{\infty}
\le \tau(M) \cdot \|\hat U[\cdot, i] - U[\cdot, i]\|_{\infty}.\]
If the $\KDE_f$ subroutine guarantees that $\|\hat U[\ell] - U[\ell]\|_{\infty} \le (3n\tau(M))^{-1} \cdot \|\bo\|_{1} = (3\tau(M))^{-1}$ for all $\ell \in \{1, \cdots, D\}$, then the entry-wise difference between $\hat W$ and $W$ is bounded by $1/3$.

The $D$ calls to the subroutine then take in total $D \cdot T(n, m, (3n\tau(M))^{-1})$ time. It takes in addition $D \cdot O(nm) = n^{1 + o(1)}$ operations to construct the vectors and $O(D^{\omega})$ time for step 2.\footnote{We always assume $f(x)$ can be exactly computed in constant time for any $x \in [0, 1]$.}
\end{proof}

We now combine the reduction above and the hardness result of BCP in Theorem \ref{thm:mainbcp}. The hardness of the KDE problem follows.
\begin{theorem} \label{thm:kdehardness}
Assuming SETH, for every $q > 0$, there exists $C \ge 0$ such that $\KDE_f(n, m, (3n\tau(M))^{-1})$
cannot be solved in time $O(n^{2-q})$ for any constant $q>0$. if either of the following holds: (1) $m > C \log n, D = m$ or (2) $m > C^{\log^* n}, D > m^{C^{\log^* n} \cdot (\log n)/m}$. Here $M$ is the $D \times D$ counting matrix in Definition \ref{def:countingmatrix}.
\end{theorem}

\section{Schur polynomials}
\label{sec:Schur}

Let $\F$ be a field. For matrices $A, B \in \F^{m \times n}$, let $A \circ B$ denote the Hadamard product of $A$ and $B$, i.e. the $m \times n$ matrix with entries $(A \circ B)[i, j] = A[i, j] \cdot B[i, j], i\in [m], j \in [n]$. For matrix $A \in \F^{m \times n}$ we define the Hadamard powers $A^{\circ 1} = A$ and $A^{\circ (k+1)} = A^{\circ k} \circ A$ for integer $k \ge 1$.
Similarly one can define the Hadamard product and Hadamard power for vectors $u \in \F^m$. Moreover, given a vector $u \in \F^m$ and a tuple $r = (r_1, \cdots, r_n) \in \N^{n}$,
we denote by $u^{\circ r}$ the $m \times n$ matrix
\[\begin{pmatrix}
| & | & & | \\
u^{\circ r_1} & u^{\circ r_2} & \cdots & u^{\circ r_n} \\
| & | & & |
\end{pmatrix}
=\begin{pmatrix}
u_1^{r_1} & u_1^{r_2} & \cdots & u_1^{r_n} \\
\vdots & \vdots & \ddots & \vdots \\
u_m^{r_1} & u_m^{r_2} & \cdots & u_m^{r_n}
\end{pmatrix}.\]
In particular, when $m = n$, we call $u^{\circ r}$ a generalized Vandermonde matrix. This is a natural generalization of the (usual) Vandermonde matrix $u^{\circ \delta}$ associated with the tuple $\delta = (0, 1, \cdots, m-1)$.

The concept of Schur polynomials was first proposed by Cauchy and defined as the ratio of two generalized Vandermonde determinants. In what follows we denote by $\N^m_< = \{x \in \N^m: x_1 < \cdots < x_m\}$ the set of $m$-tuples composed of distinct entries in ascending order, and define $\N^m_{\le} = \{x \in \N^m: x_1 \le \cdots \le x_m\}$ similarly.
\begin{definition}[Cauchy's definition of Schur polynomials]
Let $m>0$ be an integer and $\lambda = (\lambda_1, \cdots, \lambda_m) \in \N^m_{\le}$ an integer tuple.
We define the Schur polynomial on variables $(u_1, \cdots, u_m)$ by
\[s_{\lambda}(u) = s_{\lambda}(u_1, \cdots, u_m)
=\frac{\det(u^{\circ (\lambda + \delta)})}{\det(u^{\circ \delta})}
=\frac{\det(u^{\circ (\lambda + \delta)})}{V(u)}.\]
Here $V(u) = \det(u^{\circ \delta}) = \prod_{1 \le i < j \le m}(u_j - u_i)$ is the Vandermonde determinant.
\end{definition}

It is known that Schur polynomial has an equivalent definition by Littlewood using Young tableaux.
\begin{definition}
Let $m > 0$ be an integer, and $\lambda = (\lambda_1, \cdots, \lambda_m) \in \N^m_{\le}$ an integer tuple.
A semi-standard Young tableaux (SSYT) of shape $\lambda$ on alphabet $[m]$
is a left-aligned two-dimensional rectangular array $T$ of cells,
with $\lambda_i$ cells in the $i$-th row $(i \in [m])$, from bottom to top, such that
\begin{itemize}
\item each cell in $T$ is assigned with an entry from $1, \cdots, m$;
\item entries weakly decrease in each row, from left to right;
\item entries strictly decrease in each column, from top to bottom.
\end{itemize}
Moreover, for a SSYT $T$, we define the type $\mathrm t(T) = (t_1, \cdots, t_m) \in \N^m$ of $T$,
where $t_j$ is the number of cells in $T$ assigned with entry $j \in [m]$.
\end{definition}
\begin{definition}[Littlewood's definition of Schur polynomials]
Let $m > 0$ be an integer.
For integer tuple $\lambda = (\lambda_1, \cdots, \lambda_m) \in \N_{\le}^m$,
the Schur polynomial $s_{\lambda}(u)$ can be equivalently defined by
\[s_{\lambda}(u) = \sum_{T \in \mathcal T} u^{\mathrm t(T)},\]
where $\mathcal T$ is the set of all SSYT of shape $\lambda$ on alphabet $[m]$.
\end{definition}






























\section{Direct calculation of $\tau(M)$: Gaussian kernel and $t$-Student kernel}
\subsection{Gaussian kernel}
In this section we show upper bounds for 
\[\tau(M) = \max_{b \in \R^D: b \ne 0} \frac{\|M^{-1}b\|_{\infty}}{\|b\|_{\infty}} \le D \cdot \max_{t, s \in [D]} |M^{-1}[t, s]|\]
where $M$ is the counting matrix associated with the Gaussian kernel $f(x) = e^{-Bx}$.
By standard facts in linear algebra, we can rewrite
\[|M^{-1}[t, s]| = \left|\frac{\det(M^{s-}_{t-})}{\det(M)}\right|\]
with $M_{t-}^{s-}$ the submatrix of $M$ consisting of all entries but those in row $s$ or column $t$.

For the simplicity of analysis, we study the counting matrix $M$ of dimension $(D+1)$. A slight modification of the main reduction can employ such a matrix to recover the distance matrix with a redundant row: $W[p, i] = \#[\|x^{(i)} - y^{(j)}\|_2^2 = p], i \in [n], p = \{0, 1, \cdots, D\}$.
\begin{theorem}
Let $\alpha \in \R^{D+1}$ be a fixed vector and $\beta \in \R^{D+1}$ be the identity vector.
Let $M$ be the $(D+1) \times (D+1)$ counting matrix associated with Gaussian kernel $f = e^{-Bx}$ and $\alpha, \beta$. Then there exists a vector $\alpha$ such that
\[\tau(M) \le \left(\frac{5e}{1 - e^{-B/D}}\right)^D.\]
\end{theorem}

\begin{proof}
Let $x_i = \exp(-B\alpha_i)$, then
$M = [x_i^j: i, j \in \{0, \cdots, D\}]$ is a Vandermonde matrix with
\[|\det(M)| = |\det(x^{\circ(0, \cdots, D)})| = V(x).\]
On the other hand, we observe that $M_{t-}^{s-}$ can be viewed as a generalized Vandermonde matrix. Making use of both the algebraic and combinatorial definition of Schur polynomials, we have\footnote{The Schur polynomial here in fact equals an elementary symmetric polynomial.}
\[|\det (M_{t-}^{s-})| = |\det((x^{s-})^{\circ(0, \cdots, t-1, t+1, \cdots, D)})|
= V(x^{s-}) \left( \sum_{R\subseteq \{0, \cdots, D\}\backslash\{s\}}\prod_{r\in R}x^r \right),\]
in which $x^{s-}$ denotes the vector $(x_0, \cdots, x_{s-1}, x_{s+1}, \cdots, x_D)$ and
\[\sum_{R\subseteq \{0, \cdots, D\}\backslash\{s\}}\prod_{r\in R}x^r 
\le \prod_{i \in \{0, \cdots, D\} \backslash \{s\}}(1 + x_i) \le 2^D.\]
Therefore
\[\max_{s, t} |\det (M_{t-}^{s-})|
\le 2^D \cdot \max_s V(x^{s-}) = 2^D \cdot V(x)
\left( \min_s \prod_{i \in [D]\backslash \{s\}} |x_i - x_s| \right)^{-1}. \]

We now pick $0 = \alpha_0 < \cdots < \alpha_{D-1} < \alpha_{D} = 1/D$ such that
$|x_{i+1} - x_i| = |\exp(-B\alpha_{i+1}) - \exp(-B\alpha_i)|
= \frac 1{D}(1 - e^{-B/D})$. Let $r = \frac 1{D}(1 - e^{-B/D})$, then
\begin{align*}
\prod_{i \in \{0, \cdots, D\} \backslash \{s\}} |x_i - x_s|
&= \left(\prod_{i=0}^{s-1} (s - i) r \right)
\left(\prod_{i=s+1}^{D} (i-s) r\right) \\
&= r^{D} \cdot s! (D-s)!
\ge r^{D} \cdot (\frac{D}{2e})^{D/2 \cdot 2} = (\frac{1 - e^{-B/D}}{2e})^{D}.
\end{align*}
Combining the calculations together, we have
\[\tau(M) \le D \cdot \max_{s, t}\left|\frac{\det(M^{s-}_{t-})}{\det(M)}\right| 
\le D \cdot 2^D\left( \min_s \prod_{i \in [D]\backslash \{s\}} |x_i - x_s| \right)^{-1} \le \left(\frac{5e}{1 - e^{-B/D}}\right)^D.\]
\end{proof}

The following hardness result for Gaussian KDE therefore arises from a combination of the general KDE hardness Theorem \ref{thm:kdehardness} and the bound on $\tau(M)$ above.
\begin{theorem}[Hardness of Gaussian KDE]
Let $f(x) = e^{-Bx}$ be the Gaussian kernel with $B \ge 1$.
Then assuming SETH, for every $q > 0$, there exists $C_1, C_1', C_2, C_3, C_3', C_4 \ge 0$ such that $\KDE_f(n, m, \eps)$ cannot be solved in time $O(n^{2-q})$ if either of the following holds:
\begin{enumerate}[(1)]
\item $m > C_1\log n, B < C_1'\log n, 1/\eps > (C_2m/B)^m$, or
\item $C_3\log^* n < m < C_3'(\log n)/(\log\log n)$, $\log\log(1/\eps) > (\log n)/m \cdot (\log m) \cdot C_4^{\log^*n}$.
\end{enumerate}
\end{theorem}
\begin{proof}
For regime (1), first pick appropriate $C_1'$ so that $B/m < 1$.
\[3n\tau(M) = 3n\left(\frac{5e}{1 - e^{-B/m}}\right)^m \le 3n\left(\frac{5e}{B/(em)}\right)^m
=3 \cdot 2^{C_1^{-1}m} \cdot \left(\frac {5e^2m}B\right)^m \le \left(\frac {C_2m}B\right)^m\]
for some constant $C_2 > 0$.
For regime (2), we have $B/D \le 1$ and thus
\[3n\left(\frac{5e}{1 - e^{-B/D}}\right)^D \le 3n\left(\frac{5e^2D}{B}\right)^D \le 3n(5e^2D)^D.\]
If $m < (\log n)/(\log \log n)$,
\[\log D = (\log n)/m \cdot (\log m) \cdot C^{\log^* m}
= (\log\log n)^2(1 - o(1)) \cdot C^{\log^* m} > \log \log n.\]
\[\log(3n\tau(M)) \le \log(3n) + D\log(5e^2D) < C'D\log(5e^2D)\]
For some $C' > 0$. Thus,
\[\log \log(3n\tau(M)) \le \log D + \log\log(5e^2D) \le (\log n)/m \cdot (\log m) \cdot C_4^{\log^* m}.\]

\end{proof}


\subsection{$t$-Student kernel}
For $t$-Student kernels $f(x) = 1/(1 + x^\rho)$, we prove a similar bound on $\tau(M)$. The key observation here is that both $M$ and $M_{s-}^{t-}$ are (scaled) Cauchy matrices. For vectors $a, b \in \R^n$, the Cauchy matrix associated with $a, b$ is defined by $M[i, j] = 1/(a_i + b_j), i, j \in [n]$. The following closed-form formula is known for its determinant.
\begin{theorem}[Cauchy determinant]
Let $n \ge 1$ be an integer. For $a, b \in \R^{n}$, let $M = [1/(a_i + b_j)]_{i, j \in [n]}$ be the associated Cauchy matrix. Then
\[\det(M) = \frac{\prod_{1 \le i < j \le n}(a_i - a_j)(b_i - b_j)}
{\prod_{i, j \in [n]}(a_i + b_j)}.\]
\end{theorem}
\begin{corollary}\label{cor:cauchyratio}
For $a, b \in \R^n$ with $a_i, b_j \in (0, 1), \forall i, j \in [n]$, 
and $s, t \in [n]$, we have
\begin{align*}\det\left[\frac1{1+a_ib_j}\right]_{i, j \in [n]}
&= \left(\prod_{i=1}^n \frac1{a_i}\right)\det\left[\frac1{a_i^{-1}+b_j}\right] \\
&= \left(\prod_{i=1}^n \frac1{a_i}\right)
\frac{\prod_{1 \le i < j \le n}(a_i^{-1} - a_j^{-1})(b_i - b_j)}
{\prod_{i, j \in [n]}(a_i^{-1} + b_j)} 
= \frac{\prod_{1 \le i < j \le n}(a_j - a_i)(b_i - b_j)}
{\prod_{i, j \in [n]}(1 + a_ib_j)}.
\end{align*}
\[\left|\frac{\det[(1+a_ib_j)^{-1}]_{i, j \in [n]}}
{\det[(1+a_ib_j)^{-1}]_{i \in [n]^{s-}, j \in [n]^{t-}}}\right|
= \left|\frac{\prod_{i\in [n]^{s-}}(a_i - a_s)\prod_{j\in [n]^{t-}}(b_j - b_t)}
{\prod_{i \in [n]}(1 + a_ib_t)\prod_{j \in [n]^{t-}}(1 + a_sb_j)}\right|.\]
Here $[n]^{s-} = [n] \backslash\{s\}, [n]^{t-} = [n] \backslash \{t\}$.
\end{corollary}
Before proving the bound on $\tau(M)$, we gather several simple inequalities to be used.
\begin{lemma} 
(I) Let $r \ge 1, a, b > 0$ be real numbers. Then $(a+b)^r \ge a^r + b^r$.
(II) Let $a, b \ge 0$ be integers. Then $a!b! \ge ((\lfloor\frac{a+b}2\rfloor)!)^2$. 
\end{lemma}
\begin{proof}
(I) Taking the derivative, one can show $f(x) = (1 + x)^r - x^r$ is increasing in $x$ for $x > 0$.
Thus $f(b/a) = (1 + b/a)^r - (b/a)^r \ge f(0) = 1$. The inequality follows.

(II) For $a = b = 0$, the inequality is trivial. If not, we have 
$\binom{a+b}{a} \le \binom{a+b}{\lceil(a+b)/2\rceil}$, $a!b! \ge 
(\lceil\frac{a+b}2\rceil)!(\lfloor\frac{a+b}2\rfloor)! \ge ((\lfloor\frac{a+b}2\rfloor)!)^2$.
\end{proof}
\begin{theorem}
Let $\rho \ge 1$ be a real number. Let $\alpha \in \R^D$ be a fixed vector and $\beta \in \R^D$ be the identity vector.
Let $M$ be the $D \times D$ counting matrix associated with $t$-Student kernel $f(x) = 1/(1+x^{\rho})$ and $\alpha, \beta$. Then there exists a vector $\alpha$ such that
\[\tau(M) \le (7e)^{2\rho D} .\]
\end{theorem}
\begin{proof}
The counting matrix is
\[M = [f(\alpha_i\beta_j)]_{i, j \in [D]} = 
\left[\frac{1}{1 + (\alpha_i\beta_j)^{\rho}}\right]_{i, j \in [D]}\]
For simplicity we assume $D$ is odd. Let $s, t \in [D]$. By corollary \ref{cor:cauchyratio}, we have
\[\left|\frac{\det(M)}{\det(M^{s-}_{t-})}\right|
=\left|\frac{\prod_{i\in [D]^{s-}}(\alpha_i^\rho - \alpha_s^\rho)
\prod_{j\in [D]^{t-}}(\beta_j^\rho - \beta_t^\rho)}
{\prod_{i \in [D]}(1 + \alpha_i^\rho\beta_t^\rho)
\prod_{j \in [D]^{t-}}(1 + \alpha_s^\rho \beta_j^\rho)}\right|.\]
We set $\alpha = \beta$ to be the scaled identity vector with $\alpha_i = \beta_i = i/D, \forall i \in [D]$.
(By rescaling $\alpha_i = i/D^2, \beta_i = i$ one can make $\beta$ the identity vector.)
Then
\[\prod_{i \in [D]}(1 + \alpha_i^\rho\beta_t^\rho)
\prod_{j \in [D]^{t-}}(1 + \alpha_s^\rho \beta_j^\rho) \le 2^{2D-1},\]
\begin{align*}\prod_{i\in [D]^{s-}}|\alpha_i^\rho - \alpha_s^\rho|
&=\prod_{i = 1}^{s-1}\frac{s^\rho - i^\rho}{D^{\rho}}
\prod_{i = s+1}^{D}\frac{i^\rho - s^\rho}{D^{\rho}}
\ge \prod_{i = 1}^{s-1}\frac{(s-i)^\rho}{D^{\rho}}
\prod_{i = s+1}^{D}\frac{(i-s)^\rho}{D^{\rho}} \\
&=\left(\frac{(s-1)!(D-s)!}{D^{D-1}}\right)^\rho
\ge \left(\frac{(\lfloor\frac{D-1}2\rfloor)!}{D^{D-1}}\right)^{\rho}
\ge \left(\frac{D-1}{2eD}\right)^{\rho(D-1)} \ge (3e)^{-\rho D}.
\end{align*}
Similarly,
\[\prod_{j\in [D]^{t-}}|\beta_j^\rho - \beta_t^\rho| \ge (3e)^{-\rho D}.\]
In conclusion, we have
\[\tau(M) \le D \cdot \max_{s, t}\left|\frac{\det(M^{s-}_{t-})}{\det(M)}\right|
\le D \cdot 2^{2D} \cdot (3e)^{2\rho D} \le (7e)^{2\rho D}.\]
\end{proof}

Combining the bound on $\tau(M)$ above and the general KDE hardness Theorem \ref{thm:kdehardness}, we obtain
\begin{theorem}[Hardness results of $t$-Student KDE]
Let $f(x) = 1/(1+x^\rho)$ be a $t$-Student kernel parameterized by $\rho \ge 1$ an absolute constant.
Then assuming SETH, for every $q > 0$, there exists $C_1, C_2, C_3, C_3', C_4 \ge 0$ such that $\KDE_f(n, m, \eps)$ cannot be solved in time $O(n^{2-q})$ if either of the following holds:
\begin{enumerate}[(1)]
\item $m > C_1\log n, 1/\eps > n^{C_2}$, or
\item $C_3\log^* n < m < C_3'(\log n)/(\log\log n)$, $\log\log(1/\eps) > (\log n)/m \cdot (\log m) \cdot C_4^{\log^*n}$.
\end{enumerate}
Here $C_1, C_3, C_3'$ are absolute constants while $C_2, C_4$ are dependent on $\rho$.
\end{theorem}
\begin{proof}
For regime (1), 
\[3n\tau(M) = 3n(7e)^{2\rho m} = 3n(7e)^{2\rho \cdot C_1\log n} = n^{C_2}\]
for some $C_2$ dependent on $\rho$. For regime (2), if $m < (\log n)/(\log \log n)$,
\[\log D = (\log n)/m \cdot (\log m) \cdot C^{\log^* m}
= (\log\log n)^2(1 - o(1)) \cdot C^{\log^* m} > \log \log n.\]
\[\log(3n\tau(M)) = \log(3n) + 2\log(7e)\rho D < C'_\rho D,\]
\[\log\log(3n\tau(M)) < (\log n)/m \cdot (\log m) \cdot C_4^{\log^* m}.\]
for some $C_4$ dependent on $\rho$.
\end{proof}

\section{Cauchy-Binet Expansion} \label{sec:cauchybinet}
To bound $\tau(M)$ for counting matrices that are associated with more general kernels, we establish the connection between the determinant $\det(M)$ and the Taylor coefficients of $f$ via Cauchy-Binet Expansion, based on a framework given in \cite{ACSS20}. Our improvement mainly derives from an observation that the terms in the expansion are in fact generalized Vandermonde determinants. The additional structure in the determinants allows for simplifications in analysis through Schur polynomials.

Define infinite matrices $A: [D] \times \N \rightarrow \R$,
$\tilde A: [D] \times \N \rightarrow \R$, and $B: \N \times [D] \rightarrow \R$ by
\[A[i, k] = \alpha_i^k, \quad
\tilde A[i, k] = \frac{f^{(k)}(c)}{k!}\alpha_i^k, \quad B[k, j] = \beta_j^k.\]
Suppose the Taylor series of $f(c + x)$ at $c \in [0, 1]$ is \emph{absolutely convergent} whenever $c + x \in [0, 1]$. Then for $i, j \in [D]$, assuming $c + \alpha_i \beta_j \in [0, 1]$, we have
\[M[i, j] = f(c + \alpha_i\beta_j)
= \sum_{k=0}^{\infty} \frac{f^{(k)}(c)}{k!}\alpha_i^k \beta_j^k
=\sum_{k=0}^{\infty}\tilde A[i, k]B[k, j].\]
By Cauchy-Binet formula, we expand $\det(M)$ as follows.
\begin{align}
\det(M) &= \sum_{\bn \in \N^D_<} \det \tilde A[[D]; \bn] \cdot \det B[\bn; [D]] \notag\\
&=\sum_{\bn \in \N^D_<} \left(\prod_{n\in \bn} \frac{f^{(n)}(c)}{n!}\right)
\det A[[D]; \bn] \cdot \det B[\bn; [D]]  \notag\\
&=\sum_{\bn \in \N^D_<} \left(\prod_{n\in \bn} \frac{f^{(n)}(c)}{n!}\right)
\det (\alpha^{\circ \bn})\det (\beta^{\circ \bn}) \notag\\
&=\sum_{\lambda \in \N^D_{\le}}
\left(\prod_{i=1}^D \frac{f^{(\lambda_i + \delta_i)}(c)}{(\lambda_i + \delta_i)!}\right)
\det (\alpha^{\circ (\lambda + \delta)})\det (\beta^{\circ (\lambda + \delta)})
\quad (\lambda = \bn - \delta) \notag\\
&=V(\alpha)V(\beta) \sum_{\lambda \in \N^D_{\le}}
\left(\prod_{i=1}^D \frac{f^{(\lambda_i + \delta_i)}(c)}{(\lambda_i + \delta_i)!}\right)
s_{\lambda}(\alpha)s_{\lambda}(\beta). \label{eq:CBexpansion1} 
\end{align}
The last step follows from the algebraic definition of Schur polynomials. Similarly,
\begin{equation}
\det(M_{t-}^{s-}) = V(\alpha^{s-})V(\beta^{t-}) \sum_{\lambda \in \N^{D-1}_{\le}}
\left(\prod_{i=1}^{D-1} \frac{f^{(\lambda_i + \delta_i)}(c)}{(\lambda_i + \delta_i)!}\right)
s_{\lambda}(\alpha^{s-})s_{\lambda}(\beta^{t-}). \label{eq:CBexpansion2}
\end{equation}

\subsection{Absolutely monotonic kernels}
At this point one may try bounding (\ref{eq:CBexpansion1}) and (\ref{eq:CBexpansion2}) using analytic properties of Schur polynomials. However, for a general function $f$, the arbitrary signs of Taylor coefficients largely complicate the analysis, making tight bounds out of reach.
Hence we first study the easy case where $\tau(M)$ is associated with an \emph{absolutely monotonic} kernel $f$, i.e., $f:[0, 1]\rightarrow [0, 1]$ satisfies $f^{(n)}(c) \ge 0$ for all $n \in \N$ and $c > 0$. We will shortly see that such kernels are in fact not artificial as all the positive definite kernels naturally reduce to this case.

We first use the following indentity of Schur polynomials to ``align'' (\ref{eq:CBexpansion1}) and (\ref{eq:CBexpansion2}).
\begin{proposition}
If $\lambda_1 \ne 0$, then
$s_{(\lambda_1, \cdots, \lambda_n)}(\alpha_1 = 0, \alpha_2, \cdots, \alpha_n) =0$.
If $\lambda_1 = 0$, then
$s_{(\lambda_1, \cdots, \lambda_n)}(\alpha_1 = 0, \alpha_2, \cdots, \alpha_n) =
s_{(\lambda_2, \cdots, \lambda_n)}(\alpha_2, \cdots, \alpha_n)$.
\end{proposition}
\begin{proof}
By the combinatorial definition of Schur polynomials, we have
\[s_{(\lambda_1, \cdots, \lambda_n)}(\alpha_1, \alpha_2, \cdots, \alpha_n)
=\sum_{T \in \mathcal T} \alpha^{\mathrm t(T)}\]
where $\mathcal T$ is the set of all SSYT of shape $\lambda = (\lambda_1, \cdots, \lambda_n)$
on alphabet $[n]$, and $\mathrm t(T)$ denotes the type of SSYT $T$.
Now set $\alpha_1 = 0$. If the letter 1 appears in a SSYT $T$, i.e. $\mathrm t(T)[1] > 0$, then the corresponding monomial vanishes when evaluated. Therefore the equation simplifies to
\[s_{(\lambda_1, \cdots, \lambda_n)}(\alpha_1, \alpha_2, \cdots, \alpha_n)
=\sum_{T\in \mathcal T_1} \alpha^{\mathrm t(T)}\]
where $\mathcal T_1$ is the set of all SSYT of shape $\lambda$ on alphabet $\{2, 3, \cdots, n\}$.

However, if $\lambda_1 > 0$, no sequence of length $n$ on alphabet $\{2, 3, \cdots, n\}$
satisfies the (strictly) decreasing constraint in the first column of the tableau.
In this case $s_{(\lambda_1, \cdots, \lambda_n)}(\alpha_1, \alpha_2, \cdots, \alpha_n) = 0$.
If $\lambda_1 = 0$, then $\mathcal T_1$ is essentially the set of all SSYT of shape $\lambda'
=(\lambda_2, \cdots, \lambda_n)$ over alphabet $\{2, 3, \cdots, n\}$.
By definition,
\[s_{(\lambda_1, \cdots, \lambda_n)}(\alpha_1, \alpha_2, \cdots, \alpha_n)
=\sum_{T\in \mathcal T_1} \alpha^{\mathrm t(T)}
=s_{(\lambda_2, \cdots, \lambda_n)}(\alpha_2, \cdots, \alpha_n).\]
\end{proof}
In consequence, if we fix $\alpha = (\alpha_1, \alpha_2, \cdots, \alpha_D)$
with $\alpha_1 = 0$, then
\begin{align*}
\det(M) &= V(\alpha)V(\beta) \sum_{\lambda \in \N^D_{\le}}
\left(\prod_{i=1}^D \frac{f^{(\lambda_i + \delta_i)}(c)}{(\lambda_i + \delta_i)!}\right)
s_{\lambda}(\alpha)s_{\lambda}(\beta)\\
&= V(\alpha)V(\beta) \sum_{\lambda^{1-} \in \N^{D-1}_{\le}}
\left(f(c) \prod_{i=2}^{D} \frac{f^{(\lambda_i + \delta_i)}(c)}{(\lambda_i + \delta_i)!}\right)
s_{\lambda^{1-}}(\alpha^{1-})s_{\lambda^{1-}}(\beta^{1-})\\
&= V(\alpha)V(\beta) f(c) \sum_{\lambda \in \N^{D-1}_{\le}}
\left(\prod_{i=1}^{D-1} \frac{f^{(\lambda_i + \delta_{i+1})}(c)}{(\lambda_i + \delta_{i+1})!}\right)
s_{\lambda}(\alpha^{1-})s_{\lambda}(\beta^{1-}).
\end{align*}
Here $\lambda^{1-}, \alpha^{1-}, \beta^{1-}$ are obtained by removing the first entry in the corresponding partition/vector $\lambda, \alpha, \beta$.
(In the last step we abuse the notation by renaming $\lambda^{1-}$ as $\lambda$.) Meanwhile,
\begin{align*}
\max_{s, t}\det(M_{t-}^{s-}) &\le \max_{s, t} V(\alpha^{s-})V(\beta^{t-}) \cdot
\max_{s, t} \sum_{\lambda \in \N^{D-1}_{\le}}
\left(\prod_{i=1}^{D-1} \frac{f^{(\lambda_i + \delta_i)}(c)}{(\lambda_i + \delta_i)!}\right)
s_{\lambda}(\alpha^{s-})s_{\lambda}(\beta^{t-}) \\
&= \left(\max_{s, t} V(\alpha^{s-})V(\beta^{t-})\right)
\sum_{\lambda \in \N^{D-1}_{\le}}
\left(\prod_{i=1}^{D-1} \frac{f^{(\lambda_i + \delta_i)}(c)}{(\lambda_i + \delta_i)!}\right)
s_{\lambda}(\alpha^{1-})s_{\lambda}(\beta^{1-}).
\end{align*}
In what follows, we let
\[E_\lambda = \left(\prod_{i=1}^{D-1} 
\frac{f^{(\lambda_i + \delta_{i+1})}(c)}{(\lambda_i + \delta_{i+1})!}\right)
s_{\lambda}(\alpha^{1-})s_{\lambda}(\beta^{1-}), ~~
F_\lambda = \left(\prod_{i=1}^{D-1} 
\frac{f^{(\lambda_i + \delta_i)}(c)}{(\lambda_i + \delta_i)!}\right)
s_{\lambda}(\alpha^{1-})s_{\lambda}(\beta^{1-})\]
denote the corresponding terms in the two sums. For absolutely monotonic function $f$, we have $F_\lambda > 0$ and $E_{\lambda} > 0$ for all $\lambda \in N^{D-1}_{\le}$. Thus,
\begin{equation}
\label{eq:detRatio}
\max_{s, t} \frac{\det(M^{s-}_{t-})}{\det(M)}
\le \frac 1{f(c)} \max_{s, t} \frac{V(\alpha^{s-})V(\beta^{t-})}{V(\alpha)V(\beta)}
\max_{\lambda \in \N^{D-1}_{\le}}\frac{F_\lambda}{E_\lambda}.
\end{equation}

\subsection{KDE hardness for positive definite kernels} \label{sec:pdkde}
To accommodate positive definite kernels, we slightly modify the reduction in Section \ref{sec:mainreduction}.
Let $k(x, y) = f(\|x-y\|_2^2)$ be a positive definite kernel. By Schoenberg's characterization (Theorem \ref{schoenberg}), we have $(-1)^k \cdot f^{(k)}(t) \ge 0$ for all $k \in \N, t \ge 0$. Then for the function $g(x) = f(1-x)$, it holds $g^{(k)}(t) = (-1)^k \cdot f^{(k)}(1-t) \ge 0$ for all $k \in \N, t \in [0, 1]$.
Namely $g$ is an absolutely monotonic kernel. 

Next we see how $\KDE_f$ can be related to $\tau(M_g)$. Let $M$ be the $D \times D$ counting matrix associated with $g$, and let $W \in \N^{D\times n}$ be a re-indexed distance count matrix defined by $W[p, i] = \#[j\in [n]: \|x^{(i)}-y^{(j)}\|_2^2 = D-p]$. Then the matrix product $U = M \times W$ gives 
\begin{align}
U[\ell, i] &= \sum_{p=1}^D M[\ell, p] \cdot W[p, i] \notag\\ 
&=\sum_{p=1}^D g\big(\alpha_\ell\beta_p\big) 
\sum_{j \in [n]} \ind\left[\|x^{(i)} - y^{(j)}\|_2^2 = D-p\right] \notag\\
&=\sum_{j \in [n]} g\left(c + \alpha_\ell \cdot \beta\left[D-\|x^{(i)} - y^{(j)}\|_2^2\right]\right).
\end{align}
If $\beta$ is the identity vector, then
\begin{align*}
U[\ell, i] &=\sum_{j \in [n]} g\left(c + \alpha_\ell \cdot (D-\|x^{(i)} - y^{(j)}\|_2^2)\right) \\
&=\sum_{j \in [n]} f\left(1 - c - \alpha_\ell D + \alpha_\ell \|x^{(i)} - y^{(j)}\|_2^2 \right).
\end{align*}
Constructing vectors $\tilde x_{\ell}^{(i)}, \tilde y_{\ell}^{(j)} \in \R^{m+1}$ defined by $\tilde x_{\ell}^{(i)}[k] = \sqrt{\alpha(\ell)} x^{(i)}[k]$ if $k \in [m]$ and $\tilde x_{\ell}^{(i)}[m+1] = \sqrt {1-c-\alpha_{\ell}D}$;
$\tilde y_{\ell}^{(j)}[k] = \sqrt{\alpha(\ell)} y^{(j)}[k]$ if $k \in [m]$ and $\tilde y_{\ell}^{(j)}[m+1] = 0$,
we have
\[U[\ell, i] = \sum_{j \in [n]}f(\|\tilde x_{\ell}^{(i)} - \tilde y_\ell^{(j)}\|_2^2).\]
In this way we show that the new product $U$ can also be computed using KDE subroutines, and the reduction proceeds the same way as in Section \ref{sec:mainreduction} to recover the (re-indexed) distance count matrix.
By similar analysis, we have the following hardness result relating the complexity of $\KDE_f$ to $\tau(M_g)$.
\begin{theorem} \label{thm:pdkdehardness}
Let $f: [0, 1]\rightarrow [0, 1]$ be a function and $g(x) = f(1-x)$.
Then assuming SETH, for every $q > 0$, there exists $C \ge 0$ such that $\KDE_f(n, m, (3n\tau(M))^{-1})$
cannot be solved in time $O(n^{2-q})$ for any constant $q>0$. if either of the following holds: (1) $m > C \log n, D = m$ or (2) $m > C^{\log^* n}, D > m^{C^{\log^* n} \cdot (\log n)/m}$. Here $M$ is the $D \times D$ counting matrix associated with function $g$.
\end{theorem}


\subsection{Hardness of Rational Quadratic kernel}
For $f, \alpha, \beta$ of certain forms, we give explicit bounds on the two terms in (\ref{eq:detRatio}). Regarding the ratio of Vandermonde determinants, we focus on $V(x^{s-})/V(x)$ for scaled identity $x$ defined by $x_p = p/D$.
\begin{proposition}
Let $\rho \ge 1$ be a real number and $x \in [0, 1]^D$ be the vector defined by $x_i = i/D$.
Then for $s \in [D]$ we have \[\frac{V(x^{s-})}{V(x)} = \prod_{i\in [D]\backslash \{s\}}|x_i - x_s|^{-1}
\le (3e)^{D}.\]
\end{proposition}
\begin{proof}
\begin{align*}
\prod_{i\in [D]\backslash \{s\}}|x_i - x_s| 
&=\prod_{i=1}^{s-1} \frac{s - i}{D}
\cdot \prod_{i=s+1}^D \frac{i - s}{D} \\
&= \frac{(s-1)!(D-s)!}{D^{D-1}}
\ge \frac{((\frac{D-1}2)!)^2}{D^{D-1}}
\ge \left(\frac{D-1}{2eD}\right)^{(D-1)}
\ge (3e)^{-D}.
\end{align*}
\end{proof}
Hence for vectors $\alpha$ defined by $\alpha_\ell = \ell/D, \ell\in [D]$ and $\beta$ defined by $\beta_p = p/D, p \in [D]$, we have
\[\max_{s, t} \frac{V(\alpha^{s-})V(\beta^{t-})}{V(\alpha)V(\beta)} \le (3e)^{2D}.\]

For the term involving Taylor coefficients, we fix a real number $\sigma \ge 1$ and focus on the absolutely monotonic function $f(x) = (2-x)^{-\sigma}$. By calculation, for $k \in \N$,
$f^{(k)}(x) = (2-x)^{-(\sigma + k)} \prod_{i=0}^{k-1}(\sigma + i)$.
Then for $c = 0$, we have
\[\frac{F_{\lambda}}{E_{\lambda}}
=\left(\prod_{i=1}^{D-1} 
\frac{f^{(\lambda_i + \delta_i)}(0)}{(\lambda_i + \delta_i)!}\right)
\bigg/\left(\prod_{i=1}^{D-1} 
\frac{f^{(\lambda_i + \delta_{i+1})}(0)}{(\lambda_i + \delta_{i+1})!}\right)
=\prod_{i=1}^{D-1}\left(2 \cdot \frac{1 + \lambda_i + \delta_i}{\sigma + \lambda_i + \delta_i}\right)
\le 2^{D-1}.\]
Combining the calculations, we obtain the following bound on $\tau(M)$.
\begin{theorem}
Let $\sigma \ge 1$ be a real number. 
Let $f:[0, 1] \rightarrow [0, 1]$ be the function $f(x) = (2-x)^{-\sigma}$, and let vectors $\alpha, \beta \in \R^D$ be defined by $\alpha_\ell = \ell/D, \beta_p = p/D$. Then the $D \times D$ counting matrix $M$ associated with $f, \alpha, \beta$ has
\[\tau(M) \le D \cdot 2^\sigma \cdot (3e)^{2D} \cdot 2^{D-1} \le 2^{\sigma}(7e)^{2D}.\]
\end{theorem}

Combining the bound on $\tau(M)$ associated with $f(x) = (1 + (1-x))^{-\sigma}$ with the hardness of positive definite KDE Theorem \ref{thm:pdkdehardness}, we obtain
\begin{theorem}[Hardness of Rational Quadratic KDE]
Let $f(x) = 1/(1+x)^\sigma$ be a Rational Quadratic kernel parameterized by $\sigma \ge 1$ an absolute constant.
Then assuming SETH, for every $q > 0$, there exists $C_1, C_2, C_3, C_3', C_4 \ge 0$ such that $\KDE_f(n, m, \eps)$ cannot be solved in time $O(n^{2-q})$ for any constant $q>0$. if either of the following holds:
\begin{enumerate}[(1)]
\item $m > C_1\log n, 1/\eps > n^{C_2}$, or
\item $C_3\log^* n < m < C_3'(\log n)/(\log\log n), \log\log(1/\eps) > (\log n)/m \cdot (\log m) \cdot C_4^{\log^*n}$.
\end{enumerate}
Here $C_1, C_3, C_3'$ are absolute constants while $C_2, C_4$ are dependent on $\sigma$.
\end{theorem}
\begin{proof}
For regime (1), 
\[3n\tau(M) = 3n\cdot 2^{\sigma}(7e)^{2m} = n^{C_2}\]
for some $C_2$ dependent on $\rho$. For regime (2), if $m < (\log n)/(\log \log n)$,
\[\log D = (\log n)/m \cdot (\log m) \cdot C^{\log^* m} > \log \log n.\]
\[\log(3n\tau(M)) = \log(3n) + \sigma + 2\log(7e)\rho D < C'_\sigma D,\]
\[\log\log(3n\tau(M)) < (\log n)/m \cdot (\log m) \cdot C_4^{\log^* m}.\]
for some $C_4$ dependent on $\rho$.
\end{proof}























































































































































































































